% !PRINTED 23
\weberpage{23}
\thispagestyle{empty}
\begin{center}
{\large\gesperrt{Erster Abschnitt.}}

\vspace{0.8em}
{\Large\gesperrt{Rationale Functionen.}}

\vspace{0.8em}
\rule{0.16\linewidth}{0.4pt}

\vspace{0.8em}
{\large \S.\ 1.}

\vspace{0.5em}
{\large\gesperrt{Ganze Functionen.}}
\end{center}

Nächster Gegenstand der Betrachtung sind \gesperrt{ganze rationale}
oder auch kurz \gesperrt{ganze Functionen einer Veränderlichen}. Wir
verstehen darunter Ausdrücke von folgender Form:
\begin{equation}
 f(x)=a_0x^n+a_1x^{n-1}+a_2x^{n-2}+\cdots+a_{n-1}x+a_n,
\tag{1}
\end{equation}
worin $n$ ein ganzzahliger Exponent, der \gesperrt{Grad} der Function $f(x)$ ist.
Der Grad ist eine \gesperrt{natürliche Zahl}. Bisweilen ist aber auch
nützlich, von ganzen rationalen Functionen
0\textsuperscript{ten} Grades zu sprechen, worunter eine Constante
verstanden wird. $x$ heisst die \gesperrt{Veränderliche},
$a_0,a_1,a_2,\ldots,a_{n-1},a_n$ die \gesperrt{Coëfficienten}. Sowohl
$x$ als $a_0,a_1,\ldots,a_n$ sind Symbole für unbestimmte Grössen, mit
denen nach den Regeln der Buchstabenrechnung verfahren wird, für die
auch unter Umständen bestimmte Zahlwerthe gesetzt werden können (vgl.
die Einleitung).

Wenn die Function $f(x)$ in der Weise wie in (1) geschrieben ist, so
nennen wir sie nach \gesperrt{absteigenden Potenzen von $x$ geordnet}.
Die Summanden können in jeder beliebigen anderen Reihenfolge
angeordnet, also z.~B. auch nach aufsteigenden Potenzen von $x$
geordnet sein:
\[
 f(x)=a_n+a_{n-1}x+\cdots+a_1x^{n-1}+a_0x^n.
\]
Durch Addition (Subtraction) und Multiplication ganzer rationaler
Functionen entstehen wieder ganze rationale Functionen. Die
Vorschriften der Buchstabenrechnung geben unmittelbar die
Bildungsgesetze dieser neuen Functionen.

Bei der Addition ist der Coëfficient irgend einer Potenz $x^\nu$ in
der Summe gleich der Summe der Coëfficienten von $x^\nu$ in den
einzelnen Summanden. Der Grad der entstandenen Function ist
%CONT>>
